\chapter{Why Learn Ancient Greek?}

Ancient Greek is the language of Homer, of the classical dramatists, of the
New Testament, and of the philosophers whose ideas still shape how we think.
This textbook takes you from the alphabet to reading real sentences of
Classical Greek in a few chapters of steady, structured work.

\section{What you will be able to do}

By the end of this book you will be able to read adapted passages of Attic
prose, use a dictionary and grammar confidently, and understand how Greek
words are built so that long vocabulary lists stop being intimidating.

\begin{figure}[ht]
  \centering
  \includegraphics[width=0.86\linewidth]{assets/parthenon-west.jpg}
  \caption{The Parthenon on the Acropolis of Athens, seen from the west.
  The classical city produced the dialect this book teaches.}
\end{figure}

\section{How this book works}

Each chapter introduces one grammar topic, one slice of vocabulary, and one
\emph{reading} section that uses everything you have learned so far. Do the
short exercises in writing; Greek is learned with the hand as much as with
the eye.\footnote{Ten minutes of written practice a day beats an hour of
passive reading once a week.}

\begin{figure}[ht]
  \centering
  \includegraphics[width=0.6\linewidth]{assets/homer-bm-1825.jpg}
  \caption{A marble bust of Homer, Roman copy after a Hellenistic original;
  photographed in the British Museum. Tradition names Homer as the poet of
  the \emph{Iliad} and the \emph{Odyssey}.}
\end{figure}

\PSEIndex{Ancient Greek}
\PSEIndex{Attic dialect}
